% ══════════════════════════════════════════════════════════════════
%  Section 5 — Empirical Evidence
% ══════════════════════════════════════════════════════════════════

\section{Empirical Evidence}\label{sec:evidence}

Theory is beautiful.  Evidence is convincing.  This section presents a
real session---\emph{software\_quality\_discipline}---that ran to
completion under NEOS dynamics.  Every number reported below was
produced by the field engine; none was hand-tuned.

% ─────────────────────────────────────────────────────────────────
\subsection{Session Overview}

Table~\ref{tab:session-metrics} summarises the session.

\begin{table}[htbp]
\centering
\caption{Summary metrics for the \emph{software\_quality\_discipline}
  session.}
\label{tab:session-metrics}
\begin{tabular}{@{}lr@{}}
\toprule
\thead{Metric} & \thead{Value} \\
\midrule
Parameters                   & $\lambda{=}0.04$, $\alpha{=}0.35$, $\tau{=}0.35$, $\sigma{=}0.40$ \\
Patterns injected            & 69 \\
Patterns surviving           & 61 (59 ceiling, 2 asymptotic) \\
Absorptions                  & 7 (57\% self-referential) \\
Expulsions                   & 1 (Singleton, cycle 52) \\
Cycles to ground state       & 52 \\
Final coherence              & 0.993 \\
Fixed-point events           & 4 ($\Psi(\text{field}) = \text{field}$) \\
Field constants              & $R_{\text{pt}} = 0.88$, $r_{\text{SOLID}} = 0.94$ \\
\bottomrule
\end{tabular}
\end{table}

\subsubsection*{Session Timeline}

Table~\ref{tab:timeline} presents the complete chronological progression
of the session across seven injection waves and 52~cycles.

\begin{table}[htbp]
\centering
\caption{Chronological session timeline: injection waves, cycle runs,
  and coherence milestones.}
\label{tab:timeline}
\begin{tabular}{@{}cllcl@{}}
\toprule
\thead{Wave} & \thead{Patterns} & \thead{Cycles} & \thead{$C$} & \thead{Key Event} \\
\midrule
1 & p001--p006 (Review) & 1--3 & $0.71$ & \texttt{defensive\_quality} attractor \\
2 & p007--p014 (Breadth) & 4--8 & $0.79$ & \texttt{holistic\_quality} attractor \\
3 & p015--p022 (Refactoring) & 9--11 & $0.80$ & Attractor $\to$ MATURE \\
4 & p023--p031 (Testing) & 12--18 & $0.91$ & Functor $F$ discovered; mass saturation \\
5 & p032--p040 (OOP/SOLID) & 19--30 & $0.955$ & 3 absorptions; SOLID Decagon \\
--- & p044 (Singleton) & 30 & $0.941$ & Immune response; 4 inhibitors \\
6 & p046--p055 (GoF) & 34--43 & $0.984$ & Great Convergence; Golden Pair \\
7 & p056--p062 (Reuse) & 44--49 & $0.990$ & 4 self-referential absorptions \\
--- & Collapse & 50--52 & $0.993$ & Singleton expelled; ground state \\
\bottomrule
\end{tabular}
\end{table}

\subsubsection*{Wave~1: Rapid Coherence from Six Patterns}

The first six injections (p001--p006) formed the \emph{defensive\_quality}
attractor in only three cycles.  Table~\ref{tab:wave1-coherence}
shows the resonance tightening: mean pairwise~$R$ rose from 0.72
to 0.78 over three cycles, while coherence jumped from
$C = 0.42$ to $C = 0.71$.

\begin{table}[htbp]
\centering
\caption{Coherence progression during Wave~1 (6 patterns, 3 cycles).}
\label{tab:wave1-coherence}
\begin{tabular}{@{}cccl@{}}
\toprule
\thead{Cycle} & \thead{$C$} & \thead{$\bar{R}$} & \thead{Event} \\
\midrule
1 & 0.42 & 0.72 & All 6 amplified, no isolation \\
2 & 0.58 & 0.75 & Cross-cluster links forming \\
3 & 0.71 & 0.78 & Attractor \emph{defensive\_quality} emerged \\
\bottomrule
\end{tabular}
\end{table}

\noindent
The peak pairwise resonance was $R = 0.91$ between
\texttt{@security\_no\_injection} and \texttt{@input\_validation}---the
strongest coupling in the entire first wave.  The field reached HIGH
coherence from only 6~patterns in 3~cycles, establishing the
resonance infrastructure that all subsequent injections built upon.

\subsubsection*{SOLID Rediscovery from First Principles}

During Waves~2--3 (p007--p030), the field received readability, naming,
DRY, encapsulation, and architecture patterns.  An \emph{architecture
quartet} emerged spontaneously \emph{before} the SOLID principles were
explicitly injected---the field rediscovered SOLID from first principles.
Once the SOLID patterns were injected, each principle paired with a GoF
technique partner (Table~\ref{tab:solid-resonance}), forming the
\textbf{SOLID Decagon} ($\alpha_1$) with field constant $R_{\text{pt}} = 0.88$
and selection function $r = 0.94$.

\subsubsection*{Waves~3--4: Refactoring and the Testing Mirror}

Wave~3 (p015--p022) introduced refactoring patterns, with
\texttt{@behavior\_preservation} $\leftrightarrow$
\texttt{@tests\_before\_refactoring} achieving $R = 0.94$---the
strongest bond in the entire session.  Wave~4 (p023--p031) injected
nine testing patterns, triggering a mass saturation event at cycle~14
where seven patterns crossed the ceiling simultaneously.  Five
production$\leftrightarrow$testing echoes emerged, seeding the
functor~$F$ that would later be recognised during eigendecomposition.
After Wave~4, coherence reached $C = 0.91$ with a six-layer topology.

\subsubsection*{Wave~5: OOP, SOLID, and the First Absorptions}

The OOP wave (p032--p040) triggered the session's first three
absorptions.  SRP was absorbed into \texttt{@single\_responsibility}
($R = 0.98$), ISP into \texttt{@api\_minimal} ($R = 0.96$), and
the Law of Demeter decomposed across three existing patterns
with 94\% semantic coverage.  S, I, and~D had been
\emph{rediscovered} before SOLID was named---the field found the
principles from first principles.  After the OOP wave completed,
coherence reached $C = 0.955$, setting the stage for the Singleton
injection.

\subsubsection*{Waves~6--7: The Great Convergence and Final Absorptions}

Ten GoF patterns (Wave~6, p046--p055) produced the ``Great
Convergence''---14 patterns reaching ceiling in 10~cycles.
\texttt{@strategy} and \texttt{@decorator} (the ``Golden Pair'')
converged simultaneously, each resonating with all five SOLID
principles.  Coherence rose from 0.962 to 0.984.

Wave~7 (p056--p062) triggered four consecutive self-referential
absorptions, including the climactic absorption~\#7 where DIP
was applied to itself.  The absorption rate accelerated across
waves: 5.0\% $\to$ 10.0\% $\to$ 13.3\% $\to$ 14.3\%,
confirming that the field's semantic space was saturating.
Coherence reached 0.990 before the final collapse.

\noindent
The session's most dramatic moment occurred at cycle~31, when
Singleton---the first adversarial pattern---was injected:

\begin{terminalbox}[Singleton Injection --- Cycle 30]
\begin{lstlisting}[style=neosshell]
nf inject "singleton_controlled_global_access" 0.65
[INJECT] singleton_controlled_global_access (s: 0.65)
  p044 | LOWEST INITIAL ACTIVATION IN FIELD HISTORY

  STRONG resonances:   0  (first pattern with ZERO)
  MODERATE resonances: 3  (allocation, factory, abs_factory)
  WEAK resonances:     6  (tension with core principles)

  LATERAL INHIBITION DETECTED:
    @test_isolation         -> singleton  INHIBIT (-0.12)
    @determinism_no_flaky   -> singleton  INHIBIT (-0.08)
    @dependency_direction   -> singleton  INHIBIT (-0.06)
    @favor_composition      -> singleton  INHIBIT (-0.05)

  Net force: -0.014/cycle  (NEGATIVE -- first in field)
  Coherence: 0.955 -> 0.941  (LARGEST SINGLE-INJECTION DROP)
\end{lstlisting}
\end{terminalbox}

% ─────────────────────────────────────────────────────────────────
\subsection{Five Eigenvectors of Software Quality}

The session collapsed to five independent axes---eigenvectors of the
resonance matrix---that together account for 100\% of the surviving
variance (Table~\ref{tab:eigenvectors}).  Crucially, these axes were
\emph{not} programmed into the system; they emerged from the resonance
dynamics alone.

\begin{table}[htbp]
\centering
\caption{Eigenvectors discovered by the
  \emph{software\_quality\_discipline} session.}
\label{tab:eigenvectors}
\begin{tabular}{@{}clcl@{}}
\toprule
\thead{\#} & \thead{Eigenvector} & \thead{Variance} & \thead{Diagnostic Question} \\
\midrule
$\lambda_1$ & Meaning $\leftrightarrow$ Mechanism
  & 34.2\% & Am I coupling to WHAT or HOW? \\
$\lambda_2$ & Principle $\leftrightarrow$ Technique
  & 22.7\% & Do I understand WHY before HOW? \\
$\lambda_3$ & Production $\leftrightarrow$ Verification
  & 18.1\% & Can I prove it works? \\
$\lambda_4$ & Restraint $\leftrightarrow$ Generalisation
  & 14.3\% & Is this abstraction earned? \\
$\lambda_5$ & Class $\leftrightarrow$ System
  & 10.7\% & Does this hold at every scale? \\
\bottomrule
\end{tabular}
\end{table}

% ── pgfplots: Eigenvector bar chart ──────────────────────────────
\begin{figure}[htbp]
\centering
\begin{tikzpicture}
\begin{axis}[
  ybar,
  bar width=18pt,
  width=0.75\textwidth,
  height=6cm,
  ylabel={Explained Variance (\%)},
  symbolic x coords={$\lambda_1$,$\lambda_2$,$\lambda_3$,
                     $\lambda_4$,$\lambda_5$},
  xtick=data,
  ymin=0, ymax=42,
  nodes near coords,
  nodes near coords align={vertical},
  every node near coord/.append style={font=\small},
  enlarge x limits=0.2,
  grid=major,
  grid style={gray!20},
  fill=neosblue!25,
  draw=neosblue!50,
]
\addplot coordinates {
  ($\lambda_1$, 34.2)
  ($\lambda_2$, 22.7)
  ($\lambda_3$, 18.1)
  ($\lambda_4$, 14.3)
  ($\lambda_5$, 10.7)
};
\end{axis}
\end{tikzpicture}
\caption{Variance explained by each eigenvector.  The dominant axis
  ($\lambda_1$) separates \emph{meaning} from \emph{mechanism}.}
\label{fig:eigenvector-bars}
\end{figure}

% ─────────────────────────────────────────────────────────────────
\subsection{Universal Invariant}

Across all 52 cycles and four fixed-point events, a single invariant
emerged as the deepest attractor in the field:

\begin{calloutbox}
$\Psi$: ``What something \textbf{means} persists; how it
\textbf{works} changes.'' --- The Universal Invariant, the deepest
attractor in the field.
\end{calloutbox}

\noindent
Every other attractor basin nests inside~$\Psi$.  It is the ground
state of the field: the configuration from which no further relaxation
is possible.

% ─────────────────────────────────────────────────────────────────
\subsection{The Verification Functor}

The session discovered a functor $F\colon \textsc{Production} \to
\textsc{Testing}$ mapping each production-code discipline to its testing
reflection.  Dependency injection serves as the natural
transformation~$\eta$ connecting the two categories.  The functor
comprises 16~morphisms, 1~splitting, and 2~dangling edges.

\noindent
Representative morphisms include:
\begin{itemize}[nosep]
  \item $F(\texttt{behavior\_preservation}) = \texttt{test\_behavior}$
  \item $F(\texttt{edge\_case\_handling}) = \texttt{edge\_cases\_boundaries}$
  \item $F(\texttt{single\_responsibility}) = \texttt{one\_assertion\_per\_concept}$
  \item $F(\texttt{error\_handling}) = \texttt{readable\_failures}$
\end{itemize}

\noindent
The keystone morphism is
$\texttt{@test\_behavior} = F(\texttt{stable\_interfaces})$.
This functor is evidence that testing is not a separate discipline
bolted onto production code but its \emph{categorical mirror},
connected by $\eta = \text{DI}$ as the natural transformation.

% ─────────────────────────────────────────────────────────────────
\subsection{Seven Attractor Basins}

The field settled into seven distinct attractor
basins (Table~\ref{tab:basins}), organised in a depth hierarchy
consistent with the attractor theory of
Hopfield~\cite{hopfield1982neural}.

\begin{table}[htbp]
\centering
\caption{Attractor basins of the
  \emph{software\_quality\_discipline} field.}
\label{tab:basins}
\begin{tabular}{@{}cllcp{5.2cm}@{}}
\toprule
\thead{Basin} & \thead{Name} & \thead{Depth} & \thead{Patterns} & \thead{Role} \\
\midrule
$\Psi$   & Universal Attractor & ground & all
  & Meaning $>$ Mechanism \\
$\alpha_1$ & SOLID Decagon & primary & 10
  & 5 principles $\times$ 5 techniques, $R = 0.88$ \\
$\alpha_2$ & Verification Mirror & primary & 13
  & Functor $F$, $\eta = \text{DI}$, 16 mappings \\
$\alpha_3$ & Craft Basin & secondary & 12
  & 12 GoF patterns (Singleton expelled) \\
$\alpha_4$ & Reuse Protocol & secondary & 9
  & Judge $\to$ Count $\to$ Extract $\to$ Share $\to$ Parameterise \\
$\alpha_5$ & Guard Basin & tertiary & 3
  & Fail fast, defend boundaries \\
$\alpha_6$ & Model Basin & tertiary & 2
  & Value objects $\leftrightarrow$ Entities \\
$\alpha_7$ & Optimize Basin & tertiary & 2
  & Performance asymptotes \\
\bottomrule
\end{tabular}
\end{table}

\noindent
The 61~surviving patterns distribute across 51~basin-memberships (overlap
factor~1.84): some patterns bridge multiple basins, acting as
connective tissue in the field's topology.

% ─────────────────────────────────────────────────────────────────
\subsection{SOLID--Technique Resonance}

Within the SOLID Decagon ($\alpha_1$), each SOLID principle
spontaneously paired with exactly one GoF technique partner
(Table~\ref{tab:solid-resonance}).  All pairings except Liskov--Composite
reached the maximum observed resonance of $R = 0.88$.

\begin{table}[htbp]
\centering
\caption{Resonance strengths between SOLID principles and their
  emergent technique partners.}
\label{tab:solid-resonance}
\begin{tabular}{@{}llc@{}}
\toprule
\thead{SOLID Principle} & \thead{Technique Partner} & \thead{$R$} \\
\midrule
S --- Single Responsibility & Strategy Pattern    & 0.88 \\
O --- Open/Closed           & Decorator Pattern   & 0.88 \\
L --- Liskov Substitution   & Composite Pattern   & 0.80 \\
I --- Interface Segregation & Facade Pattern      & 0.88 \\
D --- Dependency Inversion  & Dependency Injection & 0.88 \\
\bottomrule
\end{tabular}
\end{table}

% ─────────────────────────────────────────────────────────────────
\subsection{The Field Equation}

The five eigenvectors, weighted by their explained variance and
anchored to the universal invariant~$\Psi$, compose into a single
closed-form quality field:

\begin{eqbox}
\begin{equation}\label{eq:quality-field}
Q(x) = \Psi \!\cdot\! \bigl[\,
  .342\,\text{SOLID}
  + .227\,F
  + .181\,\text{Proto}
  + .143\,\text{Simplex}
  + .107\,\text{Scale}
\,\bigr](x)
\end{equation}
\end{eqbox}

\noindent
Equation~\eqref{eq:quality-field} is predictive: given an arbitrary
code artefact~$x$, it returns a scalar quality assessment
whose dominant term is SOLID compliance, modulated by the universal
invariant.

\subsubsection*{Complete Absorption Registry}

Table~\ref{tab:absorptions} lists all seven absorptions.  Four (57\%)
were \emph{self-referential}: the field applied the very principle
it was absorbing.

\begin{table}[htbp]
\centering
\caption{All seven absorptions in the
  \emph{software\_quality\_discipline} session.  Rows marked~$\star$
  are self-referential.}
\label{tab:absorptions}
\begin{tabular}{@{}clllc@{}}
\toprule
\thead{\#} & \thead{Attempted Pattern} & \thead{Absorbed Into} & \thead{Type} & \thead{Self-Ref} \\
\midrule
1 & SRP principle     & @single\_resp (p002) & Exact duplicate & --- \\
2 & ISP principle     & @api\_minimal (p010) & Equivalent & --- \\
3 & Law of Demeter    & @encaps $\cap$ @api $\cap$ @dep\_dir & Compound & --- \\
4 & reuse\_thru\_comp & @favor\_comp (p011) & Same intent & $\star$ \\
5 & shared\_intent    & @reduce\_dup (p023) & Complement & $\star$ \\
6 & avoid\_premature  & @rule\_of\_3 $\cap$ @dont\_force & Compound & $\star$ \\
7 & coupling\_to\_abs & @dependency\_inv (p008) & Equivalent & $\star$ \\
\bottomrule
\end{tabular}
\end{table}

\noindent
The most striking is absorption~\#7---DIP applied to itself:

\begin{terminalbox}[Absorption \#7 --- DIP Applied to Itself]
\begin{lstlisting}[style=neosshell]
nf inject "coupling_to_abstraction_not_concretion" 0.86
[ABSORBED] -> @dependency_inversion (p008)

  DIP APPLIED TO ITSELF:
    The field held the ABSTRACTION (@dependency_inversion).
    The injection offered a CONCRETION (a specific rephrasing).
    The field coupled to its abstraction and rejected the
    concretion.

  Self-referential events: 4/7 (57%)
  Fixed-point property: Psi(field) = field
\end{lstlisting}
\end{terminalbox}

\noindent
The field's absorption behaviour exhibits a fixed-point property:
applying the field's own principles to new inputs reproduces the
field.  This is precisely the condition $\Psi(\text{field}) =
\text{field}$ that defines the universal invariant as the ground state.

\subsubsection*{Saturation Analysis}

The absorption rate accelerates as the field matures: 5.0\%
(patterns~1--20), 10.0\% (21--40), 13.3\% (41--55), 14.3\%
(56--69).  This monotone increase indicates that the field's
design-time semantic space is \emph{saturating}---new injections
increasingly express content already present.  In the limit, every
novel injection would be absorbed, and the field would be closed
under its own operations.

% ─────────────────────────────────────────────────────────────────
\subsection{The Singleton Expulsion}

At cycle~52, the Singleton pattern~\cite{gamma1994design} was expelled
from the field.  Its activation dropped below the threshold~$\tau$
because it could not sustain resonance with the SOLID principles:
Singleton couples to \emph{how} (global mutable state) rather than
\emph{what} (controlled instantiation).  The decay term steadily
eroded its activation while the SOLID cluster's amplification drew
energy away from it.

This is a significant result.  The system was not told that Singleton
is problematic; the dynamics \emph{discovered} it.  The expulsion is
consistent with the well-documented critiques of Singleton in the
software-engineering literature, but here it arises from first
principles---resonance failure with the field's dominant attractor.

Singleton's only allies were \texttt{@no\_unnecessary\_allocations}
and \texttt{@factory\_method}---both members of the Optimize cluster
($\alpha_7$), itself an asymptotic satellite.  Controversial patterns
find refuge only among other outsiders, a topology consistent with
the sociology of contested design patterns.

\begin{terminalbox}[Singleton Expulsion --- Cycles 50--52]
\begin{lstlisting}[style=neosshell]
[CYCLE 50]
  p060 @mixins_traits     0.968 -> 0.994 -> CEILING
  p044 @singleton         0.383 -> 0.369   margin: 0.019

[CYCLE 51]
  p044 @singleton         0.369 -> 0.355   margin: 0.005
  CRITICAL -- one cycle from expulsion

[CYCLE 52]
  p044 @singleton         0.355 -> 0.341
  0.341 < tau (0.350)
  PATTERN EXPELLED -- FIRST IN FIELD HISTORY
    Injected: cycle 31  |  i0 = 0.65
    Expelled: cycle 52  |  A  = 0.341
    Lifespan: 21 cycles |  Cause: sustained lateral inhibition
\end{lstlisting}
\end{terminalbox}

% ── pgfplots: Coherence evolution ────────────────────────────────
\begin{figure}[htbp]
\centering
\begin{tikzpicture}
\begin{axis}[
  width=0.82\textwidth,
  height=6.5cm,
  xlabel={Cycle},
  ylabel={Coherence},
  xmin=0, xmax=55,
  ymin=0, ymax=1.08,
  grid=major,
  grid style={gray!20},
  thick,
  mark options={solid, scale=0.8},
  legend pos=south east,
  legend style={font=\small},
]

% Main coherence curve
\addplot[neosblue!70!black, thick, mark=*, mark size=1.2pt,
         smooth, samples=30, domain=0:52]
  {0.993 * (1 - exp(-0.08*x))};
\addlegendentry{Coherence $C(t)$}

% Key event markers
\node[font=\scriptsize, anchor=south west, neosblue!60!black]
  at (axis cs:0.5, 0.02) {Session start};

\draw[gray, dashed, thin] (axis cs:10, 0) -- (axis cs:10, 0.55);
\node[font=\scriptsize, anchor=south west, fill=white, inner sep=1pt]
  at (axis cs:10.5, 0.42) {First cluster};

\draw[gray, dashed, thin] (axis cs:25, 0) -- (axis cs:25, 0.87);
\node[font=\scriptsize, anchor=south west, fill=white, inner sep=1pt]
  at (axis cs:25.5, 0.72) {Major attractor};

% Singleton expulsion line
\draw[red!70!black, dashed, thick]
  (axis cs:52, 0) -- (axis cs:52, 1.05);
\node[font=\scriptsize, anchor=south east, red!70!black,
      fill=white, inner sep=1pt]
  at (axis cs:51.5, 1.02) {Singleton expelled};

% Final coherence marker
\node[circle, fill=neosblue!70!black, inner sep=1.5pt] at (axis cs:52, 0.993) {};
\node[font=\scriptsize, anchor=south east, fill=white, inner sep=1pt]
  at (axis cs:51, 0.96) {$C = 0.993$};

\end{axis}
\end{tikzpicture}
\caption{Coherence evolution over 52 cycles (fitted model:
  $C(t) = 0.993(1 - e^{-0.08t})$, not raw telemetry).  The field
  approaches its ground state monotonically, with the Singleton
  expulsion occurring at the final cycle.}
\label{fig:coherence-evolution}
\end{figure}

% ── TikZ: Attractor basin landscape ─────────────────────────────
\begin{figure}[htbp]
\centering
\begin{tikzpicture}[
  basin label/.style={font=\scriptsize\sffamily, anchor=north},
]

% Energy landscape curve (hand-tuned Bezier)
\draw[thick, neosblue!60!black]
  (0, 3.0)
  .. controls (0.6, 3.0) and (0.9, 0.2) .. (1.5, 0.2)   % Psi basin
  .. controls (2.1, 0.2) and (2.3, 2.6) .. (3.0, 2.6)    % rise
  .. controls (3.4, 2.6) and (3.6, 1.2) .. (4.0, 1.2)    % alpha1
  .. controls (4.4, 1.2) and (4.6, 2.4) .. (5.0, 2.4)    % rise
  .. controls (5.3, 2.4) and (5.5, 1.4) .. (5.8, 1.4)    % alpha2
  .. controls (6.1, 1.4) and (6.3, 2.2) .. (6.6, 2.2)    % rise
  .. controls (6.9, 2.2) and (7.1, 1.8) .. (7.4, 1.8)    % alpha3
  .. controls (7.7, 1.8) and (7.9, 2.3) .. (8.2, 2.3)    % rise
  .. controls (8.4, 2.3) and (8.6, 1.9) .. (8.8, 1.9)    % alpha4
  .. controls (9.0, 1.9) and (9.2, 2.4) .. (9.5, 2.4)    % rise
  .. controls (9.7, 2.4) and (9.9, 2.1) .. (10.2, 2.1)   % alpha5
  .. controls (10.5, 2.1) and (10.7, 2.5) .. (11.0, 2.5) % rise
  .. controls (11.2, 2.5) and (11.4, 2.2) .. (11.6, 2.2) % alpha6
  .. controls (11.9, 2.2) and (12.1, 2.6) .. (12.4, 2.6) % rise
  .. controls (12.6, 2.6) and (12.8, 2.3) .. (13.0, 2.3) % alpha7
  .. controls (13.3, 2.3) and (13.6, 3.0) .. (14.0, 3.0);% end

% Basin labels
\node[basin label] at (1.5,  0.0) {$\Psi$};
\node[basin label] at (4.0,  1.0) {$\alpha_1$};
\node[basin label] at (5.8,  1.2) {$\alpha_2$};
\node[basin label] at (7.4,  1.6) {$\alpha_3$};
\node[basin label] at (8.8,  1.7) {$\alpha_4$};
\node[basin label] at (10.2, 1.9) {$\alpha_5$};
\node[basin label] at (11.6, 2.0) {$\alpha_6$};
\node[basin label] at (13.0, 2.1) {$\alpha_7$};

% Depth annotation
\draw[<->, gray, thin] (1.5, 0.2) -- (1.5, 3.0)
  node[midway, right, font=\scriptsize] {depth};

% Y-axis label
\node[rotate=90, anchor=south, font=\small\sffamily]
  at (-0.5, 1.5) {Energy};

% X-axis label
\node[font=\small\sffamily, anchor=north] at (6.3, -0.3)
  {Configuration space};

\end{tikzpicture}
\caption{Attractor basin landscape.  $\Psi$~is the deepest basin
  (ground state); shallower basins $\alpha_1$--$\alpha_7$ are nested
  within it at increasing energy levels.}
\label{fig:basin-landscape}
\end{figure}

%% ─────────────────────────────────────────────────────────────────
\subsection{Limitations and Reproducibility}\label{sec:limitations}

The evidence presented here derives from a single session on a single
domain (software quality).  This is a proof-of-concept demonstration,
not a controlled experiment.  Key limitations include: (i)~the
eigenvector structure may be domain-specific rather than universal;
(ii)~the resonance strengths lack confidence intervals or null-model
comparison; (iii)~the coherence curve (Figure~\ref{fig:coherence-evolution})
is a fitted exponential model of the form
$C(t) = 0.993(1 - e^{-0.08t})$, not raw telemetry.  Replication
across multiple domains, operators, and substrate models is required
before the field equation can be considered validated.  The full
session logs and specification files are available at
\url{https://github.com/samuelestronati/neos}.
